% page 240 du fac-simile (image p-239 du scan)
% bloc 165.1 x 237.7 mm ; marge gauche 8.0 mm
\pgc{-12.700mm}{0.000mm}{
\l{\cel{3}232}
\l{}
\l{\sou{-ind}-.\cel{2}Sufixo qua signifikas \textquotedbl{}qua meritas esar...ata\textquotedbl{}.}
\l{}
\l{\sou{indemno}. (\sou{yuro-cienco}). To quon onu donas ad ulu por kompensar}
\l{\cel{3}lua perdi, lua spensi. - DEFIS.}
\l{}
\l{\sou{indentar}. (\sou{trans.})\cel{2}(\sou{imprimarto}).\cel{2}Retropulsar (lineo tipogra-}
\l{\cel{3}fiala) aden la longeso-kadro, pozante avan olu \textquotedbl{}lakuno-\textquotedbl{}-}
\l{\cel{3}tipi.}
\l{}
\l{\sou{indexo.}\cel{1}- I. Tabelo alfabetala,fine di libro, en qua esas men-}
\l{\cel{3}cionita la vorti precipua e la propra nomi, uzita da la}
\l{\cel{3}autoro, kun indiko pri la pagino ube esas singla de oli.-}
\l{\cel{3}II. Mikra vergo ek metalo qua, per cifro-plako, indikas la}
\l{\cel{3}kloki, la kompreseso-grado, e c. - DEFIS.}
\l{}
\l{\sou{indico}. I. Signo qua trovigas ulo. - II. (\sou{algebro}). Signo dis-}
\l{\cel{3}tingiva quan onu donas a litero qua reprezentas, en kalkulo,}
\l{\cel{3}plura grandori analoga. - III. (\sou{fiziol.})\cel{1}\sou{Refrakto-indico}: to}
\l{\cel{3}quo indikas la relato da la sinuso di la incido-angulo, a la}
\l{\cel{3}sinuso di la angulo di refrakto. - DEFIS.}
\l{}
\l{\sou{indiciono}. (\sou{kronologio)}.\cel{2}Determino di la yaro di la cikli di}
\l{\cel{3}la periodo Julianala, por datizar la akti. - DEFIS.}
\l{}
\l{\sou{indieno}.\cel{2}Kotono-stofo kun imaji imprimita, qua origine fabik-}
\l{\cel{3}esis en India. - DFS.}
\l{}
\l{\sou{indiferenta}. I. Qua ne tendencas ad ulo prefere kam ad altro.-}
\l{\cel{3}II. Qua ne interesesas da persono o da kozo, prefere kam da}
\l{\cel{3}altra persono od altra kozo. - III. Qua ne esas plu interesiva}
\l{\cel{3}kam altra kozo. - DEFIRS.}
\l{}
\l{\sou{indigo}. Farbo blua quan onu extraktas de la folii e de la stipi}
\l{\cel{3}di la L.\cel{1}\sou{indigofera}. - DEFIRS.}
\l{}
\l{\sou{indignar}. (\sou{netrans., pri, pro, kontre}). Etiko-revoltar pro la}
\l{\cel{3}spekto di konduto tre nekonvenanta. - DEFIRS.}
\l{}
\l{\sou{indijar}. (\sou{trans.})\cel{2}Ne-havar ulo quon onu bezonas. - EFIS.}
\l{}
\l{\sou{indijena}.\cel{2}Qua naskis en la lando quan lu habitas. - EFIS.}
\l{}
\l{\sou{indikar}. (\sou{trans., ad)}\cel{2}Montrar ube esas (ulu, ulo). - EFIS.}
\l{}
\l{\sou{indikativo.}\cel{3}(\sou{gram.})\cel{2}Modo di la verbo, qua expresas la ago ye}
\l{\cel{3}maniero absoluta. - DEFIS.}
\l{}
\l{\sou{indikatoro}.\cel{2}(\sou{tekn.}) Aparato qua indikas (la rapideso-grado,}
\l{\cel{3}la direcioni, la situeso di la bifurk-aguli fervoyala, la}
\l{\cel{3}aquo-quanto di lakdiero, e c.) - DEFI.}
\l{}
\l{\sou{individuo}.\cel{2}Ento qua esas unajo distingebla ek speco, genero,}
\l{\cel{3}irga membro ek societo o socio de homi. - DEFIRS.}
}
